% Source: https://artofproblemsolving.com/community/c3h2148159_remainder
% Stable identity: aops-topic-2148159-q01
Let $N$ be such a number. The numbers that leave remainder $1$ on division by $5$ and remainder $2$ on division by $7$ are those congruent to $16$ modulo $35$ (check: $16=3\cdot5+1=2\cdot7+2$; by the Chinese remainder theorem the solution is unique modulo $35$). So $N=16+35t$ for an integer $t$. Modulo $11$,
\[16+35t\equiv5+2t\equiv5\pmod{11}\iff 2t\equiv0\pmod{11}\iff t\equiv0\pmod{11}.\]
Hence $N\equiv16\pmod{385}$. The numbers $16+385s$ are $16, 401, 786, 1171,\dots$, so the largest three-digit one is $\boxed{786}$.

Check: $786=157\cdot5+1=112\cdot7+2=71\cdot11+5$.
